
\begin{table}[htbp]
\centering
\caption{Correlations With Human Trait Scores at $B=16$ and in Dense References}
\label{tab:rq3-human}
\small
\setlength{\tabcolsep}{4pt}
\begin{tabular}{@{}lccc@{}}
\toprule
Estimator &
\shortstack{$B=16$\\$r$ [95\% CI]} &
\shortstack{Dense reference\\$r$ [95\% CI]} &
\shortstack{Difference\\$\Delta r$ [95\% CI]} \\
\midrule
B1 & .625 [.603, .645] & .651 [.628, .674] & .027 [.022, .031] \\
B2 & .629 [.609, .649] & .652 [.628, .674] & .022 [.018, .027] \\
F1 & .641 [.618, .663] & .654 [.631, .676] & .013 [.010, .015] \\
F2 & .630 [.608, .652] & .652 [.629, .674] & .022 [.018, .025] \\
F3 & .629 [.607, .651] & .651 [.628, .673] & .022 [.019, .025] \\
\bottomrule
\end{tabular}
\par\smallskip
\begin{minipage}{\linewidth}
\footnotesize
\textit{Note.} Values are Pearson correlations with human trait scores,
averaged using Fisher's $z$ transformation. The $B=16$ estimates also
average across random, p50, and Pollitt selection. Dense
references are specific to each estimator. Differences are dense minus
$B=16$ correlations. The 95\% CIs reflect variation across 16 prompts;
intervals for differences use paired prompt-level observations.
B1 = simpler Bayesian; B2 = full Bayesian;
F1 = conventional frequentist; F2 = alpha-adjusted frequentist;
F3 = bootstrap-corrected F2.
\end{minipage}
\end{table}
